% Purpose:
% Define the compute-budgeted support-to-query adaptation protocol precisely.
%
% Must define:
% - support/demo pairs vs hidden query/test pairs;
% - per-task adaptation using support only;
% - final hidden-query evaluation against true held-out targets;
% - compute budgets and benchmark tiers/subsets;
% - metrics: support similarity, query similarity, query exact success,
%   final-grid match, submitted-query rate, and practicality metrics when relevant.
%
% Non-negotiable clarity points:
% - hidden-query scoring uses true held-out targets;
% - support metrics and query metrics are distinct;
% - submitted-query rate is not success;
% - diagnostic settings must be labeled distinctly from the canonical benchmark.

We study ARC in a per-task adaptation setting with a strict separation between the examples used for adaptation and the held-out examples used for evaluation. Each task contains one or more support pairs, which expose both input and output grids, and one or more hidden query pairs, which expose only the input grid during evaluation. For a fixed task, the policy is adapted using the support pairs alone. After that adaptation budget is exhausted, the final policy is evaluated on the hidden queries and scored against the true held-out targets stored in the task buffer.

The benchmark is compute-budgeted at the task level. For each task and method, we allocate a fixed number of environment timesteps to the inner-loop adaptation procedure and compare behavior as that budget increases. The completed canonical series used in the current paper package span budgets from $10$k to $100$k timesteps, with matched reporting across no-adaptation, feedforward PPO adaptation, pretrained feedforward PPO adaptation, and recurrent PPO adaptation where available. This framing keeps the comparison focused on what additional test-time computation buys on a single task, rather than conflating task difficulty with differences in offline training scale.

The protocol distinguishes support-side fitting from hidden-query transfer at the metric level. On support pairs, we report exact success, final-grid match, and mean grid similarity after adaptation. On hidden queries, we report the same grid-comparison metrics, but always against the true held-out outputs rather than any masked training target. We also track submitted-query rate because ARC agents may terminate by explicitly submitting an answer. This diagnostic is useful, but it is not itself a success metric: a policy can submit frequently and still fail every hidden query.

Hidden-query scoring is intentionally conservative. The benchmark recomputes final-grid match and similarity from the policy's terminal working grid and mask, then compares that output against \texttt{true\_test\_output\_grids} and \texttt{true\_test\_output\_masks} stored in the JaxARC task structure. Exact success requires two conditions: the final grid must match the hidden target exactly, and the policy must actually submit that grid. We separately report final-grid match so that non-submitted but correct terminal states do not get confused with successful solved-query episodes.

In the completed canonical package, we use three public ARC task families that serve different roles. MiniARC \citep{kimPlaygroundsAbstractionReasoning2022} provides the fastest slice for rapid iteration and sanity checks. ConceptARC \citep{moskvichevConceptARCBenchmarkEvaluating2023} exposes broader variation across concept groups while keeping the task family interpretable. ARC-AGI-1 entry-level training tasks from the original ARC benchmark \citep{cholletMeasureIntelligence2019} serve as the main canonical benchmark tier because they are closer to the target ARC setting while still tractable enough for repeated budget sweeps. This tiered design lets the benchmark ask the same support-to-query question across increasing task difficulty without changing the underlying scoring logic.

The protocol also keeps canonical runs separate from diagnostics. The main benchmark story uses standard support-only adaptation followed by hidden-query evaluation under the same scoring rules across methods and budgets. More permissive settings, such as positive controls with privileged information or additional query-time sampling diagnostics, are treated as auxiliary analyses and labeled explicitly when they appear. This separation matters because the paper's central claim is about what current RL-style adaptation baselines achieve under a clean support-to-query evaluation, not about the strongest number obtainable under any easier variant of the task.
